\documentclass[11pt]{article}
\usepackage[margin=1in]{geometry}
\usepackage{amsmath,amssymb,microtype}
\usepackage[hidelinks]{hyperref}

\newcommand{\Aut}{\operatorname{Aut}}

\title{The higher-rank property-$(T)$ outcome for $\Aut(F_n)$\\
is answered by arXiv:1812.03456}
\author{Literature-status note for arXiv:1712.07167}
\date{17 August 2026}

\begin{document}
\maketitle

\begin{center}
\fbox{\parbox{0.91\textwidth}{\textbf{Status.}
Explicit later-paper answer to the substantive higher-rank outcome:
$\Aut(F_n)$ has property $(T)$ for every $n\geq5$.  The later proof uses a
stronger rank-five sum-of-squares certificate, so this note does not claim that
the bare assertion ``$\Aut(F_5)$ has property $(T)$'' is by itself a black-box
inductive hypothesis, nor that Bridson's separate quotient question is solved.}}
\end{center}

\section{The source question}

Kaluba, Nowak, and Ozawa prove in arXiv:1712.07167 that $\Aut(F_5)$ has
Kazhdan's property $(T)$ \cite{KNO}.  On PDF page 2 they ask whether this fact
can be used to deduce property $(T)$ for $\Aut(F_n)$ for every $n\geq5$, as
for higher-rank lattices.  They explain that their attempted group-theoretic
induction breaks at Bridson--Vogtmann Question 12: if the image of
$\Aut(F_n)$ in a quotient of $\Aut(F_{n+1})$ is finite, must the quotient be
finite?

Thus the passage contains two levels:
\begin{enumerate}
\item the substantive outcome that every higher-rank $\Aut(F_n)$ has
property $(T)$;
\item a stronger request for an abstract induction from the rank-five fact,
via or around the quotient question.
\end{enumerate}

\section{The later higher-rank theorem}

Kaluba, Kielak, and Nowak's follow-up arXiv:1812.03456v2 states as Theorem 1
on PDF page 2 \cite{KKN}:
\[
 \boxed{\Aut(F_n)\text{ has property }(T)\text{ for every }n\geq6.}
\]
The sentence immediately following the theorem combines it with the source's
$n=5$ result and records the corollary
\[
 \Aut(F_n)\text{ has property }(T)\qquad(n\geq5).
\]
This completely answers the practical higher-rank existence question.  The
follow-up also proves a uniform Kazhdan-radius bound for the index-two groups
$\operatorname{SAut}(F_n)$.

The mechanism is genuinely extrapolative but uses more data than abstract
property $(T)$.  The authors decompose Laplacian expressions into pieces whose
rank-$m$ sum-of-squares decompositions can be symmetrized to every larger
rank.  They reduce the infinite family to a single finite rank-five
sum-of-squares certificate.  They explicitly say that the results and methods
of the source are instrumental.

\section{Scope limitation}

Because that proof uses a particular certificate in the group algebra, it
does not establish a general theorem of the form
\[
 \Aut(F_n)\text{ has }(T)\Longrightarrow \Aut(F_{n+1})\text{ has }(T)
\]
from the property alone.  Nor does the cited follow-up claim to settle the
Bridson--Vogtmann quotient problem.  The precise disposition is therefore:
\[
 \boxed{\text{all }n\geq5:\ \text{answered affirmatively;
 bare black-box induction and quotient question: not claimed.}}
\]

\begin{thebibliography}{9}
\bibitem{KNO}
M.~Kaluba, P.~W.~Nowak, and N.~Ozawa,
\emph{$\Aut(\mathbb F_5)$ has property $(T)$},
arXiv:1712.07167 (2017).

\bibitem{KKN}
M.~Kaluba, D.~Kielak, and P.~W.~Nowak,
\emph{On property $(T)$ for $\Aut(F_n)$ and $\operatorname{SL}_n(\mathbb Z)$},
arXiv:1812.03456v2 (2021).
\end{thebibliography}

\end{document}

